\chapter{Introdução}
\label{cap:introducao}

\section{Contextualização e Definição do Problema}
\label{sec:contextualizacao}

O modal ferroviário desempenha um papel estratégico na infraestrutura logística global, destacando-se pela elevada eficiência energética e capacidade de transporte de grandes volumes de cargas e passageiros em longas distâncias. No cenário industrial contemporâneo, a confiabilidade operacional e a segurança do material rodante constituem requisitos primordiais para a competitividade e sustentabilidade desse setor \cite{smith2015rolling}.

Dentre os subsistemas mecânicos críticos de vagões e locomotivas, os rodeiros ferroviários (\textit{wheelsets}) são responsáveis diretos pela sustentação da carga, guiamento dinâmico na via permanente e transmissão dos esforços de tração e frenagem. A interface rotativa entre os eixos dos rodeiros e as estruturas dos truques é viabilizada pelas caixas de mancal (\textit{axleboxes}), no interior das quais operam rolamentos de elementos rolantes submetidos a severas solicitações operacionais \cite{randall2011rolling}. Tais componentes estão continuamente expostos a cargas radiais e axiais combinadas, oscilações térmicas, vibrações de alta frequência oriundas da interação roda-trilho e choques mecânicos decorrentes de irregularidades geométricas da via e juntas de dilatação.

Sob esse regime severo de serviço, os rolamentos estão sujeitos a mecanismos progressivos de degradação estrutural, com destaque para a fadiga por contato de rolamento (\textit{Rolling Contact Fatigue} -- RCF), desgaste abrasivo, corrosão por atrito e trincas de subsuperfície. Uma vez deflagrado um defeito pontual em componentes internos -- tais como pistas interna e externa ou nos próprios elementos rolantes --, o impacto mecânico repetitivo durante a passagem das esferas ou rolos acelera a deterioração local, provocando destacamentos superficiais (\textit{spalling}). Caso a anomalia não seja identificada e contida em seus estágios iniciais, a evolução do dano pode provocar a fragmentação catastrófica do componente, culminando no travamento do eixo em plena via e, no limite, no descarrilamento da composição ferroviária, evento de severidade crítica associado a elevados prejuízos financeiros e risco iminente à vida humana.

\section{Justificativa}
\label{sec:justificativa}

Historicamente, o setor ferroviário tem empregado estratégias de manutenção preventiva baseadas em tempo ou quilometragem acumulada para a gestão de caixas de mancal. Embora reduza a incidência de falhas em serviço, essa abordagem rígida apresenta limitações operacionais expressivas: acarreta a substituição prematura de componentes que ainda dispõem de vida útil remanescente ou, inversamente, revela-se incapaz de prevenir falhas incipientes aceleradas por condições anômalas de operação antes do intervalo de revisão programado \cite{mobley2002vibration}.

Adicionalmente, as ferrovias modernas utilizam sistemas de monitoramento wayside instalados ao longo da malha, como os detectores óptico-térmicos de caixas quentes (\textit{hotbox detectors} -- HBDs). No entanto, o princípio de atuação dos HBDs reside na detecção da radiação infravermelha emitida quando o mancal atinge patamares térmicos críticos. Sob a ótica da engenharia de manutenção, o aquecimento térmico representa um sintoma terminal de falha mecânica, decorrente do atrito metal-metal generalizado e da degradação irreversível do lubrificante. Quando um alarme de HBD é deflagrado, a janela temporal para intervenção é mínima, demandando paradas emergenciais que interrompem a circulação da malha \cite{smith2015rolling}.

Em contrapartida, a manutenção preditiva baseada na condição (\textit{Condition-Based Maintenance} -- CBM) por meio da análise de vibrações acelerométricas possibilita a identificação de defeitos mecânicos em seus estágios mais embrionários, semanas ou meses antes do surgimento de calor anômalo perceptível \cite{randall2011rolling}. O impacto de um elemento rolante sobre uma microfissura gera pulsos transitórios de alta frequência que se propagam pela estrutura do mancal, permitindo que ferramentas avançadas de Processamento Digital de Sinais (PDS) e modelos de Aprendizado de Máquina (\textit{Machine Learning}) diagnostiquem o tipo e a severidade do dano com precisão quantitativa.

A validação preliminar de algoritmos inteligentes de diagnóstico requer bancos de dados experimentais controlados e reprodutíveis. Nesse contexto, o dataset de referência da \textit{Case Western Reserve University} (CWRU) consagrou-se como o padrão científico mais adotado na literatura internacional para o desenvolvimento e comparação de modelos preditivos \cite{smith2015rolling}. A aplicação desses dados viabiliza a consolidação de uma metodologia robusta que poderá, subsequentemente, ser transferida e adaptada aos sinais monitorados em rodeiros ferroviários reais.

O desenvolvimento deste trabalho insere-se no núcleo interdisciplinar da Engenharia Mecatrônica, sintetizando conceitos de dinâmica de sistemas mecânicos, instrumentação e sensoriamento eletrônico de vibrações, e computação avançada voltada à modelagem preditiva.

\section{Objetivos}
\label{sec:objetivos}

\subsection{Objetivo Geral}

Desenvolver, implementar e validar um \textit{framework} computacional em linguagem Python para processamento digital de sinais de vibração e classificação multiclasse do tipo e severidade de falhas mecânicas em rolamentos, combinando análise estocástica no domínio do tempo, análise espectral no domínio da frequência e algoritmos de Aprendizado de Máquina aplicados ao banco de dados experimental da CWRU.

\subsection{Objetivos Específicos}

Para a consecução do objetivo geral, estabelecem-se as seguintes metas específicas:
\begin{enumerate}
    \item Desenvolver um módulo de ingestão e pré-processamento de séries temporais brutas adquiridas a \SI{12}{\kilo\hertz}, implementando a segmentação em janelas discretas com rastreabilidade de linhagem de dados (\textit{data lineage});
    \item Extrair descritores estatísticos no domínio do tempo (RMS, Curtose, Fator de Crista, Desvio Padrão, Assimetria, Pico a Pico, Fator de Forma e Fator de Impulso) correlacionados à energia e à impulsividade dos impactos mecânicos;
    \item Implementar o processamento no domínio da frequência por meio da Transformada Rápida de Fourier (FFT com janela de Hanning) e da Análise de Envelope via Transformada de Hilbert, quantificando a energia associada às frequências características de defeito (BPFO, BPFI, BSF e FTF);
    \item Implementar e otimizar classificadores supervisionados baseados em \textit{Random Forest} e \textit{Support Vector Machines} (SVM);
    \item Estruturar protocolos de validação cruzada agrupada (\textit{Stratified Group K-Fold}) para garantir a estrita prevenção de vazamento de dados (\textit{data leakage}) entre janelas do mesmo ensaio físico;
    \item Avaliar o desempenho preditivo dos modelos por meio de matrizes de confusão e métricas acadêmicas (Acurácia, Precisão, Sensibilidade e \textit{F1-Score});
    \item Propor e validar uma estratégia de votação majoritária por ensaio (\textit{file-level majority voting}) para a consolidação do diagnóstico em nível de equipamento completo.
\end{enumerate}

\section{Organização do Trabalho}
\label{sec:organizacao}

Esta monografia está estruturada em cinco capítulos:
\begin{itemize}
    \item O \textbf{\autoref{cap:introducao}} contextualiza a relevância do transporte ferroviário, delineia a problemática associada à falha de rolamentos de rodeiros, apresenta a justificativa e formaliza os objetivos;
    \item O \textbf{\autoref{cap:revisao}} estabelece a fundamentação teórica sobre a dinâmica e os modos de falha em rolamentos de esferas, detalha a matemática do processamento digital de sinais nos domínios do tempo e da frequência, e descreve a formulação dos modelos de Aprendizado de Máquina empregados;
    \item O \textbf{\autoref{cap:metodologia}} descreve detalhadamente o ambiente experimental da CWRU, a arquitetura do pipeline computacional desenvolvido em Python, a rotina de segmentação por janelas, os vetores de atributos extraídos e a metodologia de treino e validação cruzada;
    \item O \textbf{\autoref{cap:resultados}} apresenta e discute os resultados obtidos, confrontando o desempenho dos classificadores, analisando a relevância dos atributos físicos para o diagnóstico e detalhando o comportamento das matrizes de confusão e da votação majoritária por ensaio;
    \item O \textbf{\autoref{cap:conclusao}} sintetiza as principais conclusões alcançadas, destaca as contribuições do trabalho e sugere propostas para pesquisas e desdobramentos futuros.
\end{itemize}
